% 附录 A Installation Procedure（原书 A-542 … A-546）
\biappendix{安装流程}{Installation Procedure}
\label{chap:install}

\section{引言}
\subsection*{Introduction}

本章描述如何下载并安装 RedHawk 程序，步骤如下：
\begin{enumerate}
  \item 从 ANSYS Customer Portal 下载软件。
  \item 执行安装。
  \item 设置 RedHawk 软件 license。
  \item 设置 RedHawk 环境。
\end{enumerate}

\section{下载 RedHawk 软件}
\subsection*{Downloading RedHawk Software}

请登录 Ansys Customer Support Center 网站（\url{https://support.ansys.com/AnsysCustomerPortal/en\_us/Downloads/Semiconductor+Products}）获取下载说明。

\section{程序安装}
\subsection*{Program Installation}

在 \texttt{<tarball\_directory>}（例如 \texttt{/disk1/ecad/}）下载软件，即拟安装 RedHawk 的位置。按以下说明操作，使用 csh/tcsh shell（或其他 shell）：
\begin{lstlisting}
% cd <tarball_directory>
% tar -xvf RedHawk_<version>.tar.xz
% (or unxz RedHawk_<version>.tar.xz )
% setenv APACHEROOT <RedHawk Installation Directory>
\end{lstlisting}

这将允许执行已安装的 RedHawk 及相关二进制文件。

安装 RedHawk 软件后，将在 \texttt{<RedHawk\_installation\_directory>} 中创建以下目录结构，以及软件 Release Notes。

\begin{center}
\begin{tabular}{ll}
\toprule
目录 & 部分内容 \\
\midrule
\texttt{license/} & ANSYS license 文件 \\
\texttt{bin/} & 包含所有 RedHawk 相关可执行文件，如 ``flow setup'' 工具 \texttt{rh\_setup.pl} 与 \texttt{gds\_setup.pl} \\
\texttt{lib/} & 支持 \texttt{bin/} 中可执行文件的共享库与文件 \\
\texttt{lmbin/} & FlexLM manager 与工具 \\
\texttt{scripts/} & 辅助 RedHawk 执行的有用 Perl 脚本 \\
\texttt{scripts/atcl} & RedHawk TCL 查询能力的有用 TCL 工具 \\
\texttt{doc/} & 软件 Release Notes 与最新 RedHawk Users' Manual \\
\bottomrule
\end{tabular}
\end{center}

\section{RedHawk 操作系统/平台支持}
\subsection*{RedHawk Operating System/Platform Support}

RedHawk 支持 RHEL（RedHat）/CentOS 与 SLES（SUSE）操作系统。

更多详情见 ANSYS Customer Portal Platform Support。（请参阅 Platform Support Strategy 与 Platform Support by Applications 文档中的 ``Semiconductor Applications'' 部分。）

\section{设置 ANSYS 半导体产品 License}
\subsection*{Setting Up the ANSYS Semiconductor Product License}

RedHawk 使用业界标准的 FlexLM licensing 方案。本节包含下载、安装与配置 Ansys Enterprise License Manager 包所需的基本信息与步骤。更详细信息强烈建议参阅 \textit{Ansys Enterprise Licensing Guide}。

\subsection{默认目录结构}
\subsubsection*{Default Directory Structure}

ANSYS License Manager 安装的默认目录结构为：

Linux：\texttt{/ansys\_inc/shared\_files/licensing}

Windows：\texttt{<OS Drive>:\textbackslash Program Files\textbackslash Ansys Inc\textbackslash Shared Files\textbackslash Licensing\textbackslash}

本包中许多文件需要特定的相对关系。在 Linux 上将 ANSYS 特定文件保持在 \texttt{shared\_files/licensing}，在 Windows 上保持在 \texttt{Shared Files\textbackslash Licensing}，可确保所需关系。可选择不同父目录，例如 \texttt{/licensing/flexnet/shared\_files/licensing}。

\subsection{Daemon 文件位置}
\subsubsection*{Daemon File Locations}

\paragraph{lmgrd}
\begin{itemize}
  \item 说明：FlexNet License Manager Daemon
  \item 默认目录：
  \begin{itemize}
    \item Linux：\texttt{/ansys\_inc/shared\_files/licensing/<platform>}
    \item Windows：\texttt{<OS Drive>:\textbackslash Program Files\textbackslash Ansys Inc\textbackslash Shared Files\textbackslash Licensing\textbackslash <platform>}
  \end{itemize}
  \item 目录要求：\texttt{lmgrd} 可位于任意目录。
\end{itemize}

\paragraph{ansyslmd}
\begin{itemize}
  \item 说明：FlexNet ANSYS Vendor Daemon
  \item 默认目录：
  \begin{itemize}
    \item Linux：\texttt{/ansys\_inc/shared\_files/licensing/<platform>}
    \item Windows：\texttt{<OS Drive>:\textbackslash Program Files\textbackslash Ansys Inc\textbackslash Shared Files\textbackslash Licensing\textbackslash <platform>}
  \end{itemize}
  \item 目录要求：\texttt{ansyslmd} 可位于任意目录。
\end{itemize}

\subsection{下载 Ansys Enterprise License Manager 包}
\subsubsection*{Downloading the Ansys Enterprise License Manager Package}

请登录 Ansys Customer Support Center 网站（\url{https://support.ansys.com/AnsysCustomerPortal/en\_us/Downloads/Semiconductor+Products}）。下载 ANSYS License Manager 安装文件需要当前技术支持协议。登录后按以下步骤操作：
\begin{itemize}
  \item 在 ANSYS Customer Portal 中，点击 Downloads $>$ Current Release。
  \item 选择安装操作系统（Windows x64 或 Linux x64）。
  \item 点击标题右侧的 + 展开 Tools 选项。
  \item 点击 ANSYS Enterprise License Manager 的 Full Package 选项。
  \item 选择所需下载目录并点击 Save。
  \item 下载完成后，使用平台标准解压工具解压包。强烈建议将文件解压到干净的空目录结构中。
\end{itemize}

\subsection{安装 Ansys Enterprise License Manager 包}
\subsubsection*{Installing the Ansys Enterprise License Manager Package}

从解压 ANSYS Enterprise License Manager 包的目录开始执行以下说明。包内容应保持提供的目录结构。
\begin{itemize}
  \item 将文件放到目标目录。
  \item 将 license 文件以 \texttt{.lic} 扩展名命名并放到以下目录：
  \begin{itemize}
    \item Linux：\texttt{/ansys\_inc/shared\_files/licensing/license\_files}
    \item Windows：\texttt{<OS Drive>:\textbackslash Program Files\textbackslash Ansys Inc\textbackslash Shared Files\textbackslash Licensing\textbackslash license\_files}
  \end{itemize}
  若有启动 license manager 的脚本，请验证 \opt{-c} 命令选项指向正确位置。
  \item 若 APACHE 的 FLEXlm（\texttt{apached}）License Manager 正在运行：
  \begin{itemize}
    \item 关闭 \texttt{apached} License Manager。
    \item 如适用，卸载 \texttt{apached} 服务。
    \item 移除任何在启动时启动 \texttt{apached} 的程序。
    \item 将任何 \texttt{apached} license 与 FlexNet options 文件复制到：
    \begin{itemize}
      \item Linux：\texttt{/ansys\_inc/shared\_files/licensing/license\_files}
      \item Windows：\texttt{<OS Drive>:\textbackslash Program Files\textbackslash Ansys Inc\textbackslash Shared Files\textbackslash Licensing\textbackslash license\_files}
    \end{itemize}
    \item 更新 license 文件使 port 号一致，并将 VENDOR 行中的 daemon 名改为 \texttt{ansyslmd}。
  \end{itemize}
  \item 启动 FlexNet：
  \begin{itemize}
    \item Linux：
\begin{lstlisting}
<install_dir>/ansys_inc/shared_files/licensing/
    linx64/lmgrd -c <path_top_license_file> -l
    <path_to_write_log_file>
\end{lstlisting}
    \item Windows：
\begin{lstlisting}
C:\Program Files\ANSYS Inc\Shared
    Files\Licensing\winx64\lmgrd -c <path_top_license_file> -l
    <path_to_write_log_file>
\end{lstlisting}
  \end{itemize}
\end{itemize}

\section{设置 RedHawk 环境}
\subsection*{Setting Up the RedHawk Environment}

\subsection{License 文件与库目录设置}
\subsubsection*{License File and Library Directory Setup}

每个运行 RedHawk 的用户都需要设置 license 文件。以下行可添加到 \texttt{.cshrc} 文件。在 csh/tcsh 环境中设置 license 文件：
\begin{lstlisting}
% setenv APACHEDA_LICENSE_FILE <absolute_path_to_license-file>
\end{lstlisting}
或：
\begin{lstlisting}
% setenv APACHEDA_LICENSE_FILE port_number@lic_server_ip_addr
% set path=($path <abs_path_to_release>/RedHawk)
\end{lstlisting}

例如：
\begin{lstlisting}
% setenv APACHEDA_LICENSE_FILE /disk1/ecad/ANSYS/license/apacheda.lic
\end{lstlisting}
或通过网络的 license 文件：
\begin{lstlisting}
% setenv APACHEDA_LICENSE_FILE 1881@129.186.1.10
% set path=($path $APACHEROOT/bin)
% rehash
\end{lstlisting}

若库目录与 RedHawk 不在同一位置，在 \texttt{.cshrc} shell 环境中设置以下环境变量：
\begin{lstlisting}
setenv LD_LIBRARY_PATH <RedHawk_installation_dir/lib>
\end{lstlisting}
或若 \texttt{\$LD\_LIBRARY\_PATH} 已存在，则改为：
\begin{lstlisting}
setenv LD_LIBRARY_PATH "$LD_LIBRARY_PATH:<RedHawk_installation_dir/lib>"
\end{lstlisting}

若在 ksh shell 环境中未执行 \texttt{source ./setup.ksh}，则在 \texttt{.kshrc} shell 环境中设置：
\begin{lstlisting}
set LD_LIBRARY_PATH=<RedHawk_installation_dir/lib>
\end{lstlisting}
或若 \texttt{\$LD\_LIBRARY\_PATH} 已存在，则执行：
\begin{lstlisting}
set LD_LIBRARY_PATH="$LD_LIBRARY_PATH:<RedHawk_installation_dir/lib>"
export LD_LIBRARY_PATH
\end{lstlisting}

\subsection{二进制设置}
\subsubsection*{Binary Setup}

RedHawk 可用二进制可在 CentOS 6.9 及 Customer Portal Platform Support Table 列出的所有兼容 Linux 版本上运行。

使用 RedHawk 前，\texttt{source .cshrc} 或刷新 \texttt{.kshrc}。或者，若在 csh 或 tcsh shell 环境中未执行 \texttt{source ./setup.csh}，则在 \texttt{.cshrc} shell 环境中设置：
\begin{lstlisting}
setenv APACHEROOT <RedHawk_installation_dir>
set PATH=($APACHEROOT/bin $path) [ if desired ]
\end{lstlisting}

若在 ksh shell 环境中未执行 \texttt{source ./setup.ksh}，则在 \texttt{.kshrc} shell 环境中设置：
\begin{lstlisting}
set APACHEROOT=<RedHawk_installation_dir>
export APACHEROOT
\end{lstlisting}

\subsection{调用}
\subsubsection*{Invocation}

为测试安装，在应用目录中执行以下命令：
\begin{lstlisting}
redhawk
\end{lstlisting}

RedHawk GUI 应出现，如图~\ref{fig:rh-initial-gui}。现在可加载数据文件并运行 RedHawk。

\figplaceholder{Figure A-1}{ANSYS RedHawk 程序初始 GUI}
\label{fig:rh-initial-gui}
